\chapter{INTRODUCTION}

\section{Background and Motivation}

The global e-commerce market reached \$5.8 trillion in 2023 and is projected to exceed
\$8 trillion by 2027 \cite{statista2024ecommerce}. Amazon, commanding approximately
37.6\% of U.S. e-commerce market share \cite{emarketer2024}, has become the de facto
competitive battleground for consumer electronics brands. Within this environment, the
ability to systematically monitor competitor pricing, inventory status, content changes,
and customer sentiment has become a critical capability for market participants.

Traditional market research methods—periodic manual audits, syndicated reports, and
point-in-time surveys—are insufficient for capturing the rapid dynamics of Amazon's
marketplace, where prices can change multiple times per day, new entrants can displace
established products within weeks, and consumer sentiment shifts in response to listing
quality updates. There is therefore a significant need for automated, continuous, and
data-driven competitive intelligence systems.

This thesis addresses this gap by designing and implementing a fully automated pipeline
that collects, stores, processes, and visualizes key competitive signals from Amazon's
consumer electronics marketplace. The system integrates web scraping, computer vision,
natural language processing, and machine learning techniques into a cohesive analytical
framework, deployed with minimal operational overhead via cloud-based scheduling.

\section{Research Objectives}

This thesis pursues four primary objectives:

\begin{enumerate}
  \item \textbf{Build a scalable data collection pipeline} capable of tracking daily
  Best Seller Rank, price, stock status, Buy Box ownership, listing content, and customer
  reviews for target products on Amazon.com.

  \item \textbf{Develop quantitative competitive metrics} including a Listing Quality
  Score (LQS), Brand Momentum Score (BMS), and K-Means price tier classification to
  enable systematic product benchmarking.

  \item \textbf{Apply NLP techniques} — specifically a pre-trained RoBERTa model — to
  extract sentiment signals from customer reviews at both the overall and aspect levels,
  producing a time-series sentiment score per product.

  \item \textbf{Deliver actionable outputs} through an automated alert engine and an
  interactive Streamlit dashboard, enabling real-time competitive monitoring without
  manual intervention.
\end{enumerate}

\section{Research Scope}

\subsection{Market Coverage}
The study focuses on \textbf{Amazon.com} (U.S. marketplace) across three consumer
electronics sub-categories selected for their competitive intensity and data richness:
\begin{itemize}
  \item \textbf{Gaming Keyboards} — mature category with strong brand differentiation
  \item \textbf{True Wireless Earbuds (TWS)} — high-growth category with frequent new entrants
  \item \textbf{Portable Chargers} — commoditized category with significant price competition
\end{itemize}

\subsection{Data Scope}
\begin{itemize}
  \item \textbf{Category tracking:} Top 50 ranked ASINs per category, collected daily
  \item \textbf{Product tracking:} 30 watchlist ASINs (10 per category), with daily
  product-level snapshots
  \item \textbf{Review tracking:} Up to 50 recent reviews per watchlist ASIN,
  collected weekly
  \item \textbf{Time period:} 15 consecutive days of data collection
\end{itemize}

\subsection{Exclusions}
This study does not analyze Amazon advertising spend data, third-party seller performance
metrics, or international marketplace variants. Price forecasting results should be
interpreted as directional trend indicators rather than precise predictions, given the
short data collection window.

\section{Research Methodology Overview}

The research adopts a \textbf{data engineering and applied machine learning} approach
consisting of five phases:

\begin{enumerate}
  \item \textbf{Data Collection:} Automated scraping via Apify's Amazon crawler actors,
  storing results in Supabase PostgreSQL with Supabase Storage for image assets.
  \item \textbf{Feature Engineering:} Computation of derived metrics (BSR velocity,
  review growth rate, discount percentage, content hash fingerprints).
  \item \textbf{Machine Learning:} K-Means clustering for price segmentation, RoBERTa
  for review sentiment, Facebook Prophet for price trend forecasting.
  \item \textbf{Alert Generation:} Rule-based detection of significant competitive events
  (price drops, stockouts, new entrants, content changes).
  \item \textbf{Visualization:} Interactive Streamlit dashboard consolidating all metrics.
\end{enumerate}

\section{Thesis Structure}

The remainder of this thesis is organized as follows:

\textbf{Chapter 2} reviews the relevant literature on e-commerce intelligence, Amazon
marketplace dynamics, web scraping methodologies, and the NLP and ML techniques applied.

\textbf{Chapter 3} describes the data collection architecture, database schema, and
collection methodology in detail.

\textbf{Chapter 4} presents the analytical methodology for each module: price tiering,
quality scoring, momentum scoring, sentiment analysis, and forecasting.

\textbf{Chapter 5} presents and discusses the empirical results across all three product
categories.

\textbf{Chapter 6} concludes with a summary of findings, research limitations, and
directions for future work.
